\documentclass[11pt]{article}

\usepackage[spanish]{babel}
\usepackage[margin=1in]{geometry}
\usepackage{amsmath, amssymb, amsfonts}
\usepackage{booktabs}
\usepackage{enumitem}
\usepackage{tcolorbox}
\usepackage{hyperref}
\usepackage{listings}
\usepackage{xcolor}
\usepackage{parskip}
\usepackage{pgfplots}

\usetikzlibrary{arrows.meta}
\pgfplotsset{compat=1.18}

\lstset{
  language=Python,
  basicstyle=\ttfamily\footnotesize,
  breaklines=true,
  columns=flexible,
  keepspaces=true,
  commentstyle=\color{gray}\itshape,
  keywordstyle=\color{blue!70!black}\bfseries,
  stringstyle=\color{green!40!black},
  frame=L,
  xleftmargin=2em,
  showstringspaces=false,
  tabsize=4
}

\title{\textbf{Laboratorio 6: Minimax \& Adversarial Search} \vspace{1em}\\
  Inteligencia Artificial - CC3085}
\author{José Antonio Mérida Castejón}
\date{\today}

\begin{document}

\maketitle

\subsection*{Task 1 - Teoría}

\textit{Responda con criterio y análisis de ingenierx. No se esperan definiciones de libro, sino que analice las
	consecuencias de las decisiones de diseño, haciendo que sus respuestas tengan justificación técnica y análisis crítico.}

\subsection*{La Explosión Combinatoria}

\textit{Durante la clase se mencionó que la complejidad de Minimax es $O(b^d)$. Con esto en mente considere un
	juego con un factor de ramificación promedio $b=30$ (como el ajedrez). Si su computadora puede analizar 1 millón
	de nodos por segundo, responda:}

\begin{itemize}[leftmargin=*, itemsep=0.8em]

	\item \textit{¿Cuánto tiempo le tomaría explorar el árbol completo hasta una profundidad de $d = 6$ turnos?
		      ¿Y hasta $d = 10$?}

	      Para calcular el tiempo de exploración, necesitamos determinar el número total de nodos a visitar. Con una
	      complejidad de $O(b^d)$ y un factor de ramificación $b=30$, tenemos:

	      \textbf{Para $d = 6$:}
	      \begin{align*}
		      \text{Nodos totales} & = 30^6 = 729,000,000 \text{ nodos}                                                   \\
		      \text{Tiempo}        & = \frac{729,000,000}{1,000,000} = 729 \text{ segundos} \approx 12.15 \text{ minutos}
	      \end{align*}

	      \textbf{Para $d = 10$:}
	      \begin{align*}
		      \text{Nodos totales} & = 30^{10} = 590,490,000,000,000 \text{ nodos}                                                    \\
		      \text{Tiempo}        & = \frac{590,490,000,000,000}{1,000,000} = 590,490,000 \text{ segundos} \approx 18.7 \text{ años}
	      \end{align*}

	      Podemos ver que el tiempo crece de manera exponencial. Mientras que para $d=6$ tendríamos que esperar
	      unos 12 minutos (tal vez aceptable), para $d=10$ necesitaríamos casi 19 años. Básicamente, cada nivel adicional
	      multiplica el tiempo por un factor de 30, lo cual hace que el problema escale de forma absolutamente inviable
	      para profundidades mayores.

	\item \textit{Basado en estos números, explique por qué Minimax puro es inviable para juegos complejos y por
		      qué la Poda Alfa-Beta no es solo una ``mejora'', sino una necesidad absoluta.}

	      Los números anteriores demuestran por qué Minimax puro es completamente inviable en la práctica.
	      Incluso con hardware moderno que procesa un millón de nodos por segundo, la explosión combinatoria nos obliga
	      a esperar años para evaluar jugadas a profundidades razonables. En un juego real de ajedrez, donde los jugadores
	      deben responder en minutos (o incluso segundos), esto simplemente no es aceptable. Además, tomando en cuenta el caso
	      de uso de los engines de ajedrez, el cual es estudiar los mejores movimientos o entender diferentes posiciones sigue
	      siendo poco práctico esperar posiblemente horas para tener un movimiento.

	      La Poda Alfa-Beta es una necesidad para hacer que estos algoritmos sean utilizables. En el mejor caso, Alfa-Beta
	      reduce la complejidad de $O(b^d)$ a $O(b^{d/2})$. Esto significa que para nuestro ejemplo con $b=30$ y $d=10$,
	      en lugar de explorar $30^{10}$ nodos, podríamos explorar aproximadamente $30^5 = 24,300,000$ nodos, reduciendo
	      el tiempo de 19 años a unos 24 segundos.

	      Adicionalmente, es importante mencionar que sin poda, estaríamos desperdiciando recursos computacionales en
	      ramas del árbol que jamás influirán en la decisión final. Si ya sabemos que una jugada nos lleva a perder,
	      no tiene sentido explorar todas las variantes de esa pérdida. Alfa-Beta aprovecha esta observación para descartar
	      subárboles completos de forma segura, sin sacrificar la optimalidad de la solución.

	      En resumen, la diferencia entre usar Minimax puro y Alfa-Beta es la diferencia entre tener un algoritmo
	      teóricamente correcto pero prácticamente inútil, y tener una herramienta que realmente puede jugar un juego
	      y se pueda utilizar para aprender o entender más sobre él.

\end{itemize}

\subsection*{Horizonte y Función de Evaluación}

\textit{Recuerde lo que se ha hablado durante la clase. Considere que en un juego real, rara vez llegamos a un nodo
	terminal (Game Over) durante la búsqueda. Cortamos la búsqueda en una profundidad $d_{max}$ y usamos una
	función heurística Eval(s). Responda:}

\begin{itemize}[leftmargin=*, itemsep=0.8em]

	\item \textit{Si su función Eval(s) es perfecta (devuelve el valor real minimax), ¿qué profundidad $d$ necesita
		      explorar para jugar perfectamente?}

	      Si nuestra función de evaluación fuese perfecta, necesitaríamos explorar únicamente $d = 1$ (un solo
	      nivel de profundidad).

	      Una función $Eval(s)$ perfecta, por definición, nos devuelve el valor minimax real de cualquier estado.
	      Básicamente, nos dice ``si juego perfectamente desde este punto, este es el resultado final que obtendré''. Si
	      tenemos esta capacidad de ``ver el futuro'', no necesitamos simularlo porque la función ya lo hizo por nosotros.

	      En términos prácticos, simplemente generamos todos los estados sucesores inmediatos (las posibles jugadas que
	      podemos hacer ahora), evaluamos cada uno con nuestra función perfecta, y elegimos la jugada que nos lleva al
	      estado con el mejor valor. No necesitamos ver más allá porque $Eval(s)$ ya consideró todas las consecuencias
	      futuras de cada decisión.

	\item \textit{Si su función Eval(s) es muy mala (casi aleatoria), ¿ayuda de algo aumentar la profundidad de
		      búsqueda? Justifique su respuesta explicando la relación entre la calidad de la heurística y la
		      profundidad de búsqueda.}

	      Cuando nuestra función $Eval(s)$ es muy mala (casi aleatoria), aumentar la profundidad de búsqueda nos permite
	      retrasar el punto en el que dependemos de esa evaluación poco confiable. Si exploramos más profundo,
	      es más probable que alcancemos estados terminales reales (victorias, derrotas, empates) donde no necesitamos
	      la heurística en absoluto porque el resultado es definitivo.

	      En conclusión, una heurística excelente nos permite jugar bien con poca profundidad (eficiente) mientras que
	      una heurística mala requiere profundidad extrema para compensar (costoso e ineficiente). La solución práctica siempre es
	      invertir en diseñar mejores heurísticas en lugar de simplemente aumentar la profundidad indiscriminadamente.

	      También, en engines de ajedrez por ejemplo lo más novedoso ha sido implementar redes neuronales para la función de evaluación. Stockfish
	      anteriormente utilizaba una heurística ``a mano'' por consenso general de granmaestros de ajedrez, pero fueron cambiadas por redes neuronales
	      (NNUE) recientemente porque otros engines empezaron a hacerlo.

\end{itemize}

\subsection*{Adversarios Irracionales}

\textit{El algoritmo Minimax asume que el oponente juega de manera óptima. Responda:}

\begin{itemize}[leftmargin=*, itemsep=0.8em]

	\item \textit{¿Qué sucede si su oponente es un novato que comete errores graves? ¿Minimax aprovechará esos
		      errores para ganar más rápido, o simplemente jugará de manera conservadora para asegurar su
		      victoria mínima garantizada?}

	      En cuanto a su planificación, \textit{Minimax} jugará de manera conservadora. Esto se debe a la suposición fundamental
	      del algoritmo: asume que el oponente siempre toma la mejor decisión posible desde su perspectiva.

	      Pensemos en un ejemplo concreto. Supongamos que Minimax tiene dos opciones:
	      \begin{itemize}
		      \item \textbf{Jugada A:} Garantiza victoria en 10 turnos contra un oponente perfecto.
		      \item \textbf{Jugada B:} Podría ganar en 3 turnos, pero solo si el oponente comete un error específico.
		            Si el oponente juega óptimamente, esta jugada lleva a empate.
	      \end{itemize}

	      Minimax \textit{siempre} elegirá la Jugada A. Desde su perspectiva, la Jugada B es arriesgada porque ``asume''
	      que el oponente jugará perfectamente y evitará la trampa, resultando en un empate. Minimax no puede ``apostar'' a
	      que el oponente cometa errores porque eso violaría su principio de minimizar la pérdida máxima.

	      En la práctica, esto significa que un agente Minimax puede parecer ``poco ambicioso'' contra jugadores débiles
	      porque no tenderá trampas. Sin embargo, sí aprovechará los errores de forma reactiva. Como el algoritmo recalcula
	      el árbol en cada turno, si el novato comete un error evidente y deja la victoria servida, el nuevo cálculo detectará
	      la ruta más corta y cambiará instantáneamente hacia la victoria rápida. Básicamente, Minimax juega como si estuviese
	      en un campeonato mundial dónde no toma riesgos innecesarios, pero es imparable cuando el oponente se equivoca.

	\item \textit{Proponga una modificación conceptual pequeña al algoritmo para que pueda explotar los errores
		      de un oponente sub-óptimo.}

	      La modificación ideal sería transformar el algoritmo a Expectiminimax. En lugar de asumir que el oponente jugará
	      perfecto y evaluar su turno tomando el valor mínimo absoluto, reemplazamos sus decisiones por nodos de probabilidad.

	      Con este cambio, el algoritmo calcula un valor esperado. Si el agente encuentra una jugada arriesgada, evaluaría el
	      porcentaje de probabilidad de que el novato caiga en la trampa versus el porcentaje de que juegue de forma óptima.
	      Si la matemática resulta a favor, el agente tomará el riesgo para intentar explotar al jugador.

	      El desafío principal de esta implementación en el mundo real es conseguir esos datos. Para que Expectiminimax funcione,
	      se necesita modelar estadísticamente al oponente. Si se tiene un historial masivo de partidas o una forma de medir
	      exactamente qué tan probable es que un humano cometa un error específico, asignar estos porcentajes al azar va a
	      arruinar al agente, haciéndolo tomar decisiones irracionales en lugar de asegurar sus victorias.

\end{itemize}

\end{document}
